\begin{scanpage}{0058}{41}
\[\left.\left.\qquad+i\frac{\sin[ka(\sin\alpha-\sin\phi)]}{\cos\frac{\alpha+\phi}{2}}\right\}\frac{e^{i(kr+\frac\pi4)}}{\sqrt{2\pi kr}}.\right\}\tag{6.16}\]
Here one must select the upper sign in front of the braces when $\phi>0$, and one must select the lower sign when $\phi<0$.

The resulting Equations (6.14) and (6.16), in contrast to Equations (6.07), satisfy the reciprocity principle. It is not difficult to establish this by verifying that Equation (6.14) is not changed with the simultaneous replacement of $\alpha$ by $\phi$ and of $\phi$ by $\alpha$, and Equation (6.16) is not changed with the replacement of $\alpha$ by $\pi+\phi$ and of $\phi$ by $\alpha-\pi$ (if $-\pi<\phi<-\frac\pi2$) and with the replacement of $\alpha$ by $\pi-\phi$ and of $\phi$ by $\pi-\alpha$ (if $\frac\pi2<\phi<\pi$).

However, the indicated equations lead to a discontinuity of the magnetic vector tangential component $H_z$ on the plane $x=0$. This is connected with the fact that, by considering the nonuniform part of the current in the vicinity of the strip's edge to be the same as on the corresponding half-plane, we actually assume the presence of currents on the entire plane containing the strip. In order to refine the resulting expressions, it is necessary to solve the problem of secondary diffraction --- that is, diffraction of the wave travelling from one edge of the strip to its other edge. In other words, it is necessary to take into account the diffraction interaction of the strip's edges. As we see, it is also necessary to take into account the secondary diffraction in the case $\alpha=\pm\frac\pi2$ when the $H_z$ component of the fringing field must equal zero.

In Chapter V, we will return to the problem of diffraction by a strip, and together with the investigation of the secondary diffraction, we will present the results of the numerical calculation based on Equations (6.07), (6.14) and (6.16).
\end{scanpage}
